%% kapa: exam=1 round=32 year=2021 period=2 form=A subject=appraisal_law subject_ko=감정평가관계법규 seq=29 number=29
%% source: Q-Net 감정평가사 기출문제 (cst00309) artlSeq=5208397 file=extracted/1008/5208397/A형/2021년도 제32회 감정평가사 1차 2교시 A형.pdf
%% conversion: kordoc 4.24.51 -> tools/gen1.py (auto); review=proofread-sonnet; flags=none
\begin{question}{29}
공간정보의 구축 및 관리 등에 관한 법령상 지목에 관한 설명으로 옳지 않은 것은?
\begin{choices}
  \item 축산업 및 낙농업을 하기 위하여 초지를 조성한 토지와 접속된 주거용 건축물의 부지는 “대”로 한다.
  \item 지목이 유원지인 토지를 지적도에 등록하는 때에는 “유”로 표기하여야 한다.
  \item 물을 상시적으로 이용하지 않고 관상수를 주로 재배하는 토지의 지목은 “전”으로 한다.
  \item 1필지가 둘 이상의 용도로 활용되는 경우에는 주된 용도에 따라 지목을 설정한다.
  \item 토지가 임시적인 용도로 사용될 때에는 지목을 변경하지 아니한다.
\end{choices}
\end{question}
\answer{2}
